\clearpage
\section*{Problem 5}
\markboth{Problem 5}{Problem 5: Part 1}
\addcontentsline{toc}{section}{Problem 5}

\textbf{Problem Statement:} Implement and analyze the Ridders' method for solving the equation $f(x) = 0$. For guidance, you might wish to read the Wikipedia page on the Ridders' method. Also feel free to consult other sources.

\subsection*{Brief Description of the method}

Given the initial bracketing interval $[x_0, x_2]$, containing \textbf{one} solution, and such that $f(x_0)$ and $f(x_2)$ have opposite signs, compute the midpoint $x_1 = (x_0 + x_2)/2$. Then define a new function $h(x) = f(x)e^{ax}$. Find the parameter $a$ that ensures $h(x_1) = \frac{1}{2}(h(x_0) + h(x_2))$. Then define $x_3$ as the $x$-intercept of the line passing through $(x_0, h(x_0))$ and $(x_2, h(x_2))$. Use $x_3$ as one of the endpoints of the next interval bracketing the solution. The other endpoint is $x_1$ if $f(x_1)f(x_3) < 0$. Otherwise, choose either $x_0$ or $x_2$ based on the requirement that the sign of $f(x)$ at the chosen point must be opposite to the sign of $f(x_3)$. Continue until the desired accuracy is reached.

\subsection*{Part 1}
\addcontentsline{toc}{subsection}{Problem 5.1}

Show that $a$ is uniquely defined and give the equation for finding it.

\noindent\rule{\textwidth}{0.4pt}

\vspace{1em}

\textbf{Solution:}

Let $x_0$ and $x_2$ be the endpoints of the bracketing interval, and $x_1 = \frac{x_0 + x_2}{2}$ the midpoint. Let $d = x_2 - x_0$.

Define $f_0 = f(x_0)$, $f_1 = f(x_1)$, $f_2 = f(x_2)$.

Ridders' method introduces a function $h(x) = f(x) e^{a x}$, and chooses $a$ so that:
\begin{align*}
    h(x_1) &= \frac{1}{2} (h(x_0) + h(x_2)) \\
    f_1 e^{a x_1} &= \frac{1}{2} \left( f_0 e^{a x_0} + f_2 e^{a x_2} \right)
\end{align*}
We know that $x_1 = x_0 + \frac{d}{2}$ and $x_2 = x_0 + d$.

Substitute and rearrange:
\begin{align*}
    f_1 e^{a(x_0 + d/2)} &= \frac{1}{2} \left( f_0 e^{a x_0} + f_2 e^{a(x_0 + d)} \right) \\
    \text{Dividing both sides by } e^{a x_0}: \\
    2 f_1 e^{a d/2} &= f_0 + f_2 e^{a d} \\
    2 f_1 e^{a d/2} - f_2 e^{a d} - f_0 &= 0
\end{align*}
Let $y = e^{a d/2}$, so $e^{a d} = y^2$.

This gives a quadratic in $y$:
\begin{align*}
    f_2 y^2 - 2 f_1 y + f_0 &= 0
\end{align*}
Solving for $y$:
\begin{align*}
    y &= \frac{2 f_1 \pm \sqrt{(2 f_1)^2 - 4 f_2 f_0}}{2 f_2} \\
    y &= \frac{f_1 \pm \sqrt{f_1^2 - f_2 f_0}}{f_2}
\end{align*}
Therefore,
\begin{align*}
    a &= \frac{2}{d} \ln\left( \frac{f_1 \pm \sqrt{f_1^2 - f_2 f_0}}{f_2} \right)
\end{align*}
We note that since $y=e^{a d/2}$, its value must be nonnegative. Since clearly $f_1^2 - f_2 f_0 > f_1^2$ ($f_2*f_0 < 0$), the sign is chosen so that $y$ is positive:
\begin{align*}
    \text{If } f_2 > 0: &\quad a = \frac{2}{d} \ln\left( \frac{f_1 + \sqrt{f_1^2 - f_2 f_0}}{f_2} \right) \\
    \text{else If } f_2 < 0: &\quad a = \frac{2}{d} \ln\left( \frac{f_1 - \sqrt{f_1^2 - f_2 f_0}}{f_2} \right)
\end{align*}
Or, more concisely:
\begin{align*}
    a &= \frac{2}{d} \ln\left( \frac{f_1 + \operatorname{sgn}(f_2) \sqrt{f_1^2 - f_2 f_0}}{f_2} \right) \\
    \text{or equivalently:}\\
\end{align*}
where $\operatorname{sgn}(f_2)$ is the sign of $f_2$.

\vspace{2em}
\noindent\rule{\textwidth}{0.4pt}
